\section{Discussion and Conclusion}

\subsection{What the results say}

The headline comparison is negative: at equal budget, a faithful AdaJEPA re-implementation, a
frozen world model and DTA-JEPA are statistically indistinguishable over 38 conditions. That
result is worth more than it looks, because the measurements that accompany it tell us where
the adaptation budget is actually being spent.

\begin{enumerate}[leftmargin=1.2em,itemsep=2pt]
\item \textbf{Recursion is the part that works.} Refinement depth one to three reduces one-step
latent error by $42\%$ in distribution and $49\%$ on unseen geometries, in a single forward
pass, with no gradient step and no extra parameters. Depth four is worse than depth three. A
world model that thinks twice in latent space is strictly better than one that thinks once, and
a world model that thinks four times is not. This is the strongest empirical claim in the paper
and it is independent of any adaptation mechanism --- with one exception that we would not have
found without running the diagnostics on both models: under a \emph{dynamics} shift the same
recursion makes error worse at every depth ($97.3 \to 130.5$ on the maze, \S\ref{app:maze}),
because the operator being iterated is itself the wrong one. The lesson is not "recurse more"
but "recurse only where the operator is trustworthy, and adapt where it is not".
\item \textbf{Allocation is bounded by calibration, not by the policy.} Our controller allocates
on the ratio of surprise to its running mean precisely because absolute surprise is
miscalibrated ($82\%$ relative calibration error, variance range $2.7$--$6.8$ against an error
range of $11.7$--$52.0$). The ratio is rank-informative ($\rho = 0.55$--$0.64$) but saturates,
so the controller ends up near its upper envelope and behaves like a uniformly stronger update.
Fixing the *scale* of the predictive variance --- not the allocation rule --- is the
prerequisite for uncertainty-allocated adaptation to deliver what the design promises.
\item \textbf{Safety mechanisms can silently become the controller.} At $\epsilon_f = 0.8$ the
anchor ball binds after every Adam step with $594.7$k fast parameters, so the effective step
size is set by the trust region rather than the learning rate. The mechanism behaved exactly as
specified --- and that is why the method stopped behaving like AdaJEPA-plus-adaptation. Any
online-adaptation method with an anchor needs the anchor radius calibrated against
$\eta\sqrt{n}$, or the constraint must be reported as part of the method.
\item \textbf{Some benchmark shifts do not test what they are named after.} A latent-distance
cost cancels photometric corruption that is applied consistently to the current and goal frame:
blur changes the one-step error by $0.02\%$, and all eight visual conditions collapse onto the
same twelve episodes. Shift suites for latent planners must corrupt the *goal* differently from
the *stream*, or measure a representation-level quantity (e.g.\ decoded or probe-based), or they
are measuring episode sampling noise.
\end{enumerate}

\subsection{Limitations}\label{sec:limits}

\paragraph{An unexercised mechanism.} The evaluation runs the episodes of a condition in
parallel from the shared pretrained checkpoint, each with its own fast-parameter copy. That
protocol exercises uncertainty allocation, episodic memory, the anchor ball and the validation
gate, but it never reaches the episode boundary at which the skill bank is written and the slow
parameters are consolidated (the reported skill-bank size is zero). The cross-episode claims of
the design are therefore implemented but \emph{not measured} here; a sequential-episode suite is
the missing experiment, and until it exists the dual-timescale title rests on the predictor's
computational split, which \emph{is} measured, rather than on cross-episode accumulation, which
is not.

\paragraph{Underpowered ablations.} Each ablation row is 84 episodes (shape) or 36 (dynamics)
against a $\pm 5$-point sampling error and a 6-point spread across rows. We report them as
inconclusive rather than selecting the ordering that flatters the method.

\paragraph{Protocol-level, not benchmark-level, replication.} Our environments are a
physics-based numpy re-implementation with $64\times64$ flat-shaded observations and
$\sim$1.3M-parameter world models trained on a few hundred trajectories per family. The shift
families, segment-based goal sampling, planner, budget, adaptation rates and success criteria
follow AdaJEPA, but the physics, the observation statistics and the encoders differ from
PushT/PointMaze with DINO-WM or ResNet encoders. Absolute success rates are not comparable
across the two papers; only comparisons under identical conditions are meaningful, which is why
every number in \S\ref{sec:results} comes from a single sweep with a fixed planner and episode
set.

\paragraph{Planner and budget.} Horizon 8 with 10 GD steps and at most 6 replans is a smaller
budget than AdaJEPA's (horizon 25, 100 GD steps, 20 replans, 100 actions). A weaker planner
compresses everyone's success rate toward the same ceiling and reduces the discriminating power
of the comparison; our absolute rates (32--44\% on manipulation, 66--70\% on navigation) are
consistent with a planner-limited regime rather than a world-model-limited one. The
mechanistic measurements (depth curves, calibration, anchor binding) do not depend on the
planner and are the parts of this study we would trust first.

\paragraph{Cheap but not free.} Adaptation costs $13\%$ of replanning time here (AdaJEPA: $5\%$),
because refinement depth multiplies the predictor cost and the OOD branch re-encodes the buffer
at every step. At $K_{\max}=4$ this is a real cost that the current results do not pay back.

\subsection{Conclusion}

We set out to make test-time adaptation of latent world models *selective*: a slow dynamics
module that should not move, a fast recursive module that should, and one surprise signal that
decides how much of everything to spend. The selective machinery is built, specified and
measured end to end. What the measurements show is that selection is not where the remaining
headroom is. Refinement depth is; the scale of the predictive variance is; and the relationship
between the anchor radius and the optimiser's natural step size is, quietly, everything. Our
contribution to the pursuit of adaptive world models is therefore as much a set of diagnostic
results --- with a fully reproduced protocol and a runnable benchmark that isolates each
mechanism --- as it is an architecture, and we report the negative comparison rather than
tuning the budget until it turns positive.

\paragraph{On STLS-JEPA.} The four repairs each did what their diagnosis predicted and the
success rate did not move. That is the useful half of a negative result: the allocation policy,
the trust region and the refinement depth are now instrumented and demonstrably behaving as
designed, which narrows the planning bottleneck to the latent \emph{geometry} rather than to the
adaptation schedule --- and the straightening experiment shows that geometry and predictive
accuracy are not the same lever, since a $64\%$ error reduction came with a halving of planning
success on the same family.
